\documentclass[french]{beamer}
% Pour le script tex2wims de Bernadette
% ======================================
%\wimsinclude{defwims_latex2wims.sty}

\def\defwims{\def\wims@dummy}  %permet l'inclusion de la commande \defwims
\def\environmentwims#1[#2][#3]#4#5#6{} %permet l'inclusion de la commande \environmentwims
\def\wimsentre#1{}   %permet l'inclusion de la commande \wimsentre  
\def\wimsinclude#1{}  %permet l'inclusion de la commande \wimsinclude  
\def\wimsnavig#1#2{} %permet l'inclusion de la commande \wimsnavig  
\def\samestyle#1{} %permet l'inclusion de la commande \samestyle
\def\typefold#1{}  %permet l'inclusion de la commande \typefold
\def\typelink#1{}   %permet l'inclusion de la commande \typelink
\def\wimsoption#1{}   %permet l'inclusion de la commande \wimsoption

% dans les deux lignes suivantes changer l'adresse (permet un lien hyperref dans
%le fichier pdf
\def\exercise#1#2{\href{http://wims.auto.u-psud.fr/wims/wims.cgi?#1&cmd=new}{WIMS : #2}}
\def\doc#1#2{\href{http://wims.auto.u-psud.fr/wims/wims.cgi?#1&cmd=new}{WIMS : #2}}
\defwims\wimsnavig#1#2{\#1{#2}}
\defwims\wimsentre#1{\wimsentre#1}

\usepackage{comment}
\wimsinclude{wims_ambali_beam.sty}

\wimsinclude{wims_ambali_style.css}


\includecomment{latexonly}
\excludecomment{wimsonly}
%%%%%%%%%%%%%%%%%%%
\title{S�quence 27 : In�quations} 
\author{Jean-Baptiste Frondas} 
\date{Pour les troisi�mes}
%%%%%%%%%%%%%%%%%%%
\usepackage{pgf,tikz}
\usetikzlibrary{arrows}
\usetikzlibrary{shapes,patterns}

\usetheme[secheader]{Madrid}

\usepackage[latin1]{inputenc}
\usepackage[T1]{fontenc}
\usepackage[upright]{fourier}
    \usepackage[scaled=0.875]{helvet}%
    \renewcommand{\ttdefault}{lmtt}%\
    \let\leq\leqslant
    \let\geq\geqslant



\definecolor{Peach}{cmyk}{0,0.50,0.70,0}
\definecolor{Purple}{cmyk}{0.45,0.86,0,0}
\definecolor{Lavender}{cmyk}{0,0.48,0,0}

\definecolor{fondregle}{cmyk}{0,0,.4,0}
\definecolor{fonddefi}{cmyk}{0.3,0,0,0}
\definecolor{fondtheo}{cmyk}{0,0.33,0.15,0}
\definecolor{shadecolor}{cmyk}{0,0,0,.2}
\definecolor{link}{cmyk}{0,0,100,0}

\beamerboxesdeclarecolorscheme{clair}{fondtheo}{fonddefi}

% D�finition
\newenvironment<>{defi}[1][]{%
\setbeamercolor{block title}{fg=black,bg=fonddefi}
\setbeamercolor{block body}{fg=black,bg=fonddefi!20!white}
    \begin{block}{D�finition{ : #1}}%
    }{%
    \end{block}}

% Th�or�me
\newenvironment<>{theo}[1][]{%
\setbeamercolor{block title}{fg=black,bg=fondtheo}
\setbeamercolor{block body}{fg=black,bg=fondtheo!20!white}
    \begin{block}{Th�or�me{ : #1}}%
    }{%
    \end{block}}

% Propri�t�s
\newenvironment<>{prop}[1][]{%
\setbeamercolor{block title}{fg=black,bg=fondtheo}
\setbeamercolor{block body}{fg=black,bg=fondtheo!20!white}
    \begin{block}{Propri�t�{ : #1}}%
    }{%
    \end{block}}

% Rappel
\newenvironment<>{rap}[1][]{%
\setbeamercolor{block title}{fg=black,bg=green}
\setbeamercolor{block body}{fg=black,bg=green!20!white}
    \begin{block}{Rappel{ : #1}}%
    }{%
    \end{block}}

% Exemple
\newenvironment<>{eg}[1][]{%
\setbeamercolor{block title}{fg=black,bg=Peach}
\setbeamercolor{block body}{fg=black,bg=Peach!20!white}
    \begin{block}{Exemple{ : #1}}%
    }{%
    \end{block}}

\newenvironment<>{egs}[1][]{%
\setbeamercolor{block title}{fg=black,bg=Peach}
\setbeamercolor{block body}{fg=black,bg=Peach!20!white}
    \begin{block}{Exemples{ : #1}}%
    }{%
    \end{block}}
% Remarque
\newenvironment<>{rmq}[1][]{%
\setbeamercolor{block title}{fg=black,bg=Purple}
\setbeamercolor{block body}{fg=black,bg=Purple!20!white}
    \begin{block}{Remarque{ : #1}}%
    }{%
    \end{block}}

	% Exercice
	\newenvironment<>{exo}[1][]{%
	\setbeamercolor{block title}{fg=black,bg=yellow}
	\setbeamercolor{block body}{fg=black,bg=yellow!20!white}
	    \begin{block}{Exercice{ : #1}}%
	    }{%
	    \end{block}}
	
	% Solution
		\newenvironment<>{sol}[1][]{%
		\setbeamercolor{block title}{fg=black,bg=yellow}
		\setbeamercolor{block body}{fg=black,bg=yellow!10!white}
		    \begin{block}{Solution{ : #1}}%
		    }{%
		    \end{block}}


\AtBeginSection[] % pour mettre la table des mati�res 
%                 % au d�but de chaque section
{ \begin{frame}<beamer> 
\frametitle{Plan de la pr�sentation} 
\tableofcontents[currentsection] 
\end{frame} }

\AtBeginSubsection[] % Do nothing for \subsection* 
{
\begin{frame}<beamer> 
%\frametitle{currentsection} 
\tableofcontents[currentsection,currentsubsection]
\end{frame}
}


\AtBeginSubsubsection[] % Do nothing for \subsection* 
{
\begin{frame}<beamer> 
%\frametitle{Outline} 
\tableofcontents[currentsection,currentsubsection]
\end{frame}
}

\usepackage{pgf,tikz}
\usetikzlibrary{arrows}

\begin{document}

\begin{latexonly}
	\begin{frame}[fragile]
		\titlepage
	\end{frame}
\end{latexonly}

\begin{wimsonly}
	\title{Sur les in�quations} 
	\author{Jean-Baptiste Frondas} 
\end{wimsonly}


\section{D�finitions} % (fold)
\label{sec:definitions}

\begin{frame}[fragile]
	\begin{defi}[In�quation]
		Une in�quation est une in�galit� dans laquelle intervient un nombre inconnu d�sign� le plus souvent par une lettre.
	\end{defi}

	\pause

	\begin{latexonly}
		\begin{egs}
			\begin{eqnarray*}
				x &<& 3 \\
				w+3 & \leq &3w+2 \\
				x^2+1 &> &3
			\end{eqnarray*}
		\end{egs}
	\end{latexonly}

\begin{wimsonly}
	\begin{wims}
		\def{integer a=random(2..12)*random(-1,1)}
		\def{integer b=random(2..12)*random(-1,1)}
		\def{integer c=random(2..12)*random(-1,1)}
		\def{integer d=random(2..12)*random(-1,1)}
		\def{text var=random(x,\a x,\a x,\a x,\a x,x^2,-2x^2)}
		\def{text var2=random(x,\b x,\b x,\b x,3x^2)}
		\def{text inf=random(>,<)}
	\end{wims}
	\begin{eg}[ \reload{<img src = "gifs/doc/etoile.gif" alt = "rechargez" width = "20" height = "20">}	]
		\begin{center}
			\( \var +\c \inf \var2 +\d \)
		\end{center}
	\end{eg}
\end{wimsonly}

	\pause

	\begin{defi}[R�soudre une in�quation]
		R�soudre une in�quation, c'est d�terminer toutes les valeurs que peut \emph{prendre} l'inconnue afin que l'in�quation soit v�rifi�e. Ces valeurs sont les solutions de l'in�quation.
	\end{defi}
\end{frame}

\begin{frame}[fragile]
	\frametitle{Tests de valeurs}
	\begin{latexonly}
		\begin{eg}
			On consid�re l'in�quation $x+3>2x-7$. La valeur 7 est-elle solution ? La valeur 12 ?\par

			On calcule chacun de deux membres avec la valeur particuli�re donn�e :\par

			Pour la valeur 7, on a : $7+3=10$ et $2\times 7-7=7$. Or $10>7$, donc 7 est une solution.
			De m�me, vous pouvez v�rifier que $0$, $-11$, $3$ sont �galement solution. il y en a une infinit�.\par

			Pour la valeur 12, on a : $12+3=15$ et $2\times 12-7=17$. Or $15<17$, donc 12 n'est pas solution.
		\end{eg}

		\begin{exo}[Laiss� au lecteur]
			On consid�re l'in�quation $4x^2-1>0$. Les valeurs suivantes sont-elles solution ? $0$, $\dfrac13$, $2$, $-3$, $1$, $-1$...
		\end{exo}
	\end{latexonly}

	\begin{wimsonly}
		\begin{wims}
			\def{integer a=random(2..12)*random(-1,1)}
			\def{integer b=random(2..12)*random(-1,1)}
			\def{integer c=random(2..12)*random(-1,1)}
			\def{integer d=random(2..12)*random(-1,1)}
			\def{integer e=random(2..12)*random(-1,1)}
			\def{integer f=random(2..12*random(-1,1))}
			\if{\e=\f}{\def{integer f=\f-3}}

			\def{function var=random(x+\c,\a *x+\c,\a *x+\c,\a *x+\c,\a x+\c,x^2+\c,-2*x^2+\c)}
			\def{function var2=random(x+\d,\b *x+\d,\b *x+\d,\b *x+\d,3*x^2+\d)}
			\def{text inf=random(>,<)}
			\def{real res1=evalue(\var,x=\e)}
			\def{real res2=evalue(\var2,x=\e)}
			\def{real toto=\res1-\res2}

			\def{real res11=evalue(\var,x=(\f))}
			\def{real res22=evalue(\var2,x=(\f))}
			\def{real toto2=(\res11)-(\res22)}

			\if{\toto<0 and \inf=< }{\def{text conc= On a bien \(\res1<\res2\) donc, \e est bien solution}}
			\if{\toto<0 and \inf=>}{\def{text conc= On a \( \res1<\res2\), donc \e n'est pas solution}}
			\if{\toto>0 and \inf=< }{\def{text conc= On a \( \res1>\res2\), donc \e n'est pas solution}}
			\if{\toto>0 and \inf=> }{\def{text conc= On a \( \res1>\res2\), donc \e est solution.}}
			\if{\toto=0}{\def{text conc=On a \(\res1=\res2\), donc \e n'est pas solution.}}


			\if{\toto2<0 and \inf=< }{\def{text conc2= On a bien \(\res11<\res22\) donc, \f est bien solution}}
			\if{\toto2<0 and \inf=>}{\def{text conc2= On a \( \res11<\res22\), donc \f n'est pas solution}}
			\if{\toto2>0 and \inf=< }{\def{text conc2= On a \( \res11>\res22\), donc \f n'est pas solution}}
			\if{\toto2>0 and \inf=> }{\def{text conc2= On a \( \res11>\res22\), donc \f est solution.}}
			\if{\toto2=0}{\def{text conc2=On a \(\res11=\res22\), donc \f n'est pas solution.}}


		\end{wims}
		\begin{eg}[ \reload{<img src = "gifs/doc/etoile.gif" alt = "rechargez" width = "20" height = "20">}	]
			On consid�re l'in�quation \(\var \inf \var2 \) . La valeur \e est-elle solution ? La valeur \f ? <br>

			On calcule chacun de deux membres avec la valeur particuli�re donn�e :<br>

			Pour la valeur \e, on a : $\var=\res1$ et $\var2=\res2$. \conc <br>


			Pour la valeur \f, on a : $\var=\res11$ et $\var2=\res22$. \conc2 <br>
		\end{eg}
	\end{wimsonly}
\end{frame}
% section definitions (end)

\section{M�thode} % (fold)
\label{sec:methode}
\begin{frame}[fragile]
	\frametitle{Comment r�soudre une in�quation ?}

	\begin{itemize}
		\item \textbf{Objectif :}  Isoler $x$ dans un membre de l'in�quation.\pause
		\item \textbf{Proc�d� :}  Transformer l'in�quation initiale en une in�quation �quivalente, \pause c'est-�-dire ayant le m�me ensemble solution.
	\end{itemize}
\end{frame}

\begin{frame}[fragile]
	\frametitle{R�gles}

	\begin{theo}[R�solution d'in�quations]
		On ne change pas l'ensemble solution d'une in�quation en : \pause
		\begin{itemize}
			\item Ajoutant un m�me nombre (positif ou n�gatif) � ses deux membres. \pause
			\item Multipliant (ou divisant) par un m�me nombre \textbf{strictement positif} les deux membres de cette in�quation. \pause
			\item Multipliant (ou divisant) par un m�me nombre \textbf{strictement n�gatif} les deux membres de cette in�quation, \textbf{� condition de changer le sens de l'in�quation.}
		\end{itemize}
	\end{theo}
\end{frame}
% section methode (end)


\section{Exemples} % (fold)
\label{sec:exemples}

\begin{frame}[fragile]
	\frametitle{Premier exemple}
	\begin{latexonly}
	On veut r�soudre l'in�quation : $3x-5 \leq 2x+1$.\pause
	\begin{eqnarray*}
		3x-5 &\leq& 2x+1 \\ \pause
		3x-5 {\textcolor{red}{+5}} & \leq& 2x+1 {\textcolor{red}{+5}} \\ \pause
		3x & \leq & 2x+6 \\ \pause
		3x {\textcolor{red}{-2x}} &\leq  & 2x+6{\textcolor{red}{-2x}} \\ \pause
		x & \leq & 6
	\end{eqnarray*}
	Conclusion : \pause Les nombres inf�rieurs ou �gaux � 6 sont solutions de l'in�quation.

	Repr�sentation graphique : \pause
	\begin{center}
		\begin{tikzpicture}[draw=black,xshift=5cm]
		 \draw (-5.5cm,0cm)--(5.5cm,0cm);
		\foreach \i in {-10,...,10}{\draw[xshift=\i*0.5 cm] (0cm,.1cm)--(0cm,-.1cm) node[yshift=-.15cm]{\i};}
		\pause
			\draw[xshift=3cm,line width=1.2pt]
			(-0.1cm,0.3cm)--(0cm,0.3cm)--(0cm,-0.3cm)--(-0.1cm,-0.3cm) node [yshift=-0.5cm] {$\displaystyle{6}$};
	  \pause
		\path[pattern= north east lines] (3cm,-0.1cm)  rectangle (5.5cm,0.25cm);
			\end{tikzpicture}
		\end{center}
		\pause
		On hachure \textbf{ce qui n'est pas solution}
	\end{latexonly}


	\begin{wimsonly}
		\begin{wims}
			\def{integer a=random(7..15)}
			\def{integer b=random(1..10)}
			\def{text s1=randitem(+,-)}
			\def{text s2=randitem(+,-)}
			 \if{\s1=+}{\def{text s1c=-}}{\def{text s1c=+}}
			 \if{\s2=+}{\def{text s2c=-}}{\def{text s2c=+}}
			\def{integer ap=random(2..\a-1)}
			\def{integer bp=random(1..10)}
			\def{text ineq=randitem(leq,geq)}
			\def{integer A=\a - \ap}
			\def{integer B=\bp \s1c \b}
			\def{rational C=\B / \A}
			\if{\ineq=leq}{\def{text graphsol=segment  -10,0,\C,0,green}}{\def{text graphsol=segment  10,0,\C,0,green}}
			\if{\ineq=leq}{\def{text hache=parallel 10.2,0.6,9.8,-0.6,-0.5,0,2*(10-\C),red}}{\def{text hache=parallel -10.2,0.6,-9.8,-0.6,0.5,0,2*abs(\C+10),red}}
			\if{\ineq=leq}{\def{text crochet = parallel \C-0.3,1,\C,1,0,-2.1,2,green}}{\def{text crochet = parallel \C+0.3,1,\C,1,0,-2.1,2,green}}

		\end{wims}
		On veut r�soudre l'in�quation : \( \a x \s1 \b \\ineq \ap x \s2 \bp \) \reload{<img src = "gifs/doc/etoile.gif" alt = "rechargez" width = "20" height = "20">}
		<br>
		On proc�de ainsi :

		\begin{eqnarray*}
			\a x \s1 \b & \\ineq & \ap x \s2 \bp \\
			\a x \s1 \b - {\textcolor{red}{\ap x}} & \\ineq & \ap x \s2 \bp - {\textcolor{red}{\ap x}} \\
			\A x \s1 \b & \\ineq & \bp \\
			\A x \s1 \b \s1c {\textcolor{red}{\b}} & \\ineq & \bp \s1c {\textcolor{red}{\b}} \\
			\A x & \\ineq & \B \\
			\frac{\A x}{{\textcolor{red}{\A}}} & \\ineq & \frac{\B}{{\textcolor{red}{\A}}} \\
			x & \\ineq & \C
		\end{eqnarray*}

	    \begin{center}
	      \draw{500,70}
	      {xrange -10,10
	      yrange -3,3
	      hline  0,0,1,0,black
	      parallel -10,-0.8,-10,0.8,1,0,20,grey
		  parallel 0,-0.8,0,0.8,1,0,2,black
		  text black,0,-1,medium,O
		  text black,1,-1,medium,1
		  segment  \C,-1,\C,1,green
		  \crochet
		  \hache
		  \graphsol
		  }
	    \end{center}
		On hachure la partie non solution (ici en rouge), et surligne (ici en vert) la partie solution. Il est important de toujours mentionner la l�gende. (Certaines personnes hachurent la partie solution, il n'y a pas de r�gle, il faut pr�ciser � chaque fois.)
	\end{wimsonly}
\end{frame}


\begin{frame}[fragile]
	\frametitle{Deuxi�me exemple}
	\begin{latexonly}
		On veut r�soudre l'in�quation : $3x+3 \leq 7x +11$. \pause
		\begin{eqnarray*}
			3x+3 & \leq & 7x +11 \\ \pause
			3x+3 {\textcolor{red}{-7x}} & \leq& 7x+11{\textcolor{red}{-7x}}  \\ \pause
			-4x+3 &\leq & 11 \\ \pause
			-4x+3{\textcolor{red}{-3}} & \leq& 11{\textcolor{red}{-3}} \\ \pause
			-4x &\leq& 8 \\ \pause
			\dfrac{-4x}{{\textcolor{red}{-4}}} & {\textcolor{red}{\geq}} & \dfrac{8}{{\textcolor{red}{-4}}} \\ \pause
			x &\geq & -2
		\end{eqnarray*}
		Repr�sentation graphique : \pause
		\begin{center}
			\begin{tikzpicture}[draw=black,xshift=5cm]
			 \draw (-5.5cm,0cm)--(5.5cm,0cm);
			\foreach \i in {-5,...,5}{\draw[xshift=\i cm] (0cm,.1cm)--(0cm,-.1cm) node[yshift=-.15cm]{\i};}
			\pause
				\draw[xshift=-2cm,line width=1.2pt]
				(0.1cm,0.3cm)--(0cm,0.3cm)--(0cm,-0.3cm)--(0.1cm,-0.3cm) node [yshift=-0.5cm] {$\displaystyle{-2}$};
		  \pause
			\path[pattern= north east lines] (-6cm,-0.1cm)  rectangle (-2cm,0.25cm);

			\end{tikzpicture}
		\end{center}
	\end{latexonly}
	\begin{wimsonly}
		\begin{wims}
			\def{integer a=random(2..7)}
			\def{integer b=random(1..10)}
			\def{text s1=randitem(+,-)}
			\def{text s2=randitem(+,-)}
			 \if{\s1=+}{\def{text s1c=-}}{\def{text s1c=+}}
			 \if{\s2=+}{\def{text s2c=-}}{\def{text s2c=+}}
			\def{integer ap=random(8..15)}
			\def{integer bp=random(1..10)}
			\def{text ineq=randitem(leq,geq)}
			\if{\ineq=leq}{\def{text ineq2=geq}}{\def{text ineq2=leq}}
			\def{integer A=\a - \ap}
			\def{integer B=\bp \s1c \b}
			\def{rational C=\B / \A}
			\if{\ineq=geq}{\def{text graphsol=segment  -10,0,\C,0,green}}{\def{text graphsol=segment  10,0,\C,0,green}}
			\if{\ineq=geq}{\def{text hache=parallel 10.2,0.6,9.8,-0.6,-0.5,0,2*(10-\C),red}}{\def{text hache=parallel -10.2,0.6,-9.8,-0.6,0.5,0,2*abs(\C+10),red}}
			\if{\ineq=geq}{\def{text crochet = parallel \C-0.3,1,\C,1,0,-2.1,2,green}}{\def{text crochet = parallel \C+0.3,1,\C,1,0,-2.1,2,green}}

		\end{wims}
		On veut r�soudre l'in�quation : \( \a x \s1 \b \\ineq \ap x \s2 \bp \) \reload{<img src = "gifs/doc/etoile.gif" alt = "rechargez" width = "20" height = "20">}
		<br>
		On proc�de ainsi :

		\begin{eqnarray*}
			\a x \s1 \b & \\ineq & \ap x \s2 \bp \\
			\a x \s1 \b - {\textcolor{red}{\ap x}} & \\ineq & \ap x \s2 \bp - {\textcolor{red}{\ap x}} \\
			\A x \s1 \b & \\ineq & \bp \\
			\A x \s1 \b \s1c {\textcolor{red}{\b}} & \\ineq & \bp \s1c {\textcolor{red}{\b}} \\
			\A x & \\ineq & \B \\
			\frac{\A x}{{\textcolor{red}{\A}}} & {\textcolor{red}{\\ineq2}} & \frac{\B}{{\textcolor{red}{\A}}} \\
			x & \\ineq2 & \C
		\end{eqnarray*}

	    \begin{center}
	      \draw{500,70}
	      {xrange -10,10
	      yrange -3,3
	      hline  0,0,1,0,black
	      parallel -10,-0.8,-10,0.8,1,0,20,grey
		  parallel 0,-0.8,0,0.8,1,0,2,black
		  text black,0,-1,medium,O
		  text black,1,-1,medium,1
		  segment  \C,-1,\C,1,green
		  \crochet
		  \hache
		  \graphsol
		  }
	    \end{center}
		On hachure la partie non solution (ici en rouge), et surligne (ici en vert) la partie solution. Il est important de toujours mentionner la l�gende. (Certaines personnes hachurent la partie solution, il n'y a pas de r�gle, il faut pr�ciser � chaque fois.)
	\end{wimsonly}
\end{frame}
% section exemples (end)

\section{Exercices} % (fold)
\label{sec:exercices}
\begin{wimsonly}
\begin{exo}
    \exercise{module=H3/algebra/eqineq.fr&+cmd=new&+worksheet=10&+mode=3&+natcoef=fractionsrel&+difficulte=2&+nocalcul=1&+nbreexo=1}{R�soudre une in�quation - M�thode}
\end{exo}
\end{wimsonly}

\begin{latexonly}
	\begin{frame}
		\frametitle{En local}
		\begin{exo}
		    \href{http://127.0.0.1/wims/wims.cgi?module=H3/algebra/eqineq.fr&+cmd=new&+worksheet=10&+mode=3&+natcoef=fractionsrel&+difficulte=2&+nocalcul=1&+nbreexo=1}{R�soudre une in�quation - M�thode}
		\end{exo}
	\end{frame}
	\begin{frame}
		\frametitle{Sur internet}
		\begin{exo}
		    \href{http://wims.auto.u-psud.fr/wims/wims.cgi?module=H3/algebra/eqineq.fr&+cmd=new&+worksheet=10&+mode=3&+natcoef=fractionsrel&+difficulte=2&+nocalcul=1&+nbreexo=1}{R�soudre une in�quation - M�thode}
		\end{exo}
	\end{frame}
\end{latexonly}
% section exercices (end)




\end{document}